% Compile twice with: pdflatex Review_Report_Class_Example.tex
% Keep ssnreview.cls and the optional logo in the same folder.
\documentclass{ssnreview}
\usepackage{amsmath,amssymb}

\projecttitle{Reasoning Failure Boundaries in ECG\\
Foundation Models Under Progressive\\
Physiological Information Loss}
\reviewnumber{REVIEW REPORT -- II}
\students{%
  Vijay K (3122235001157)\\
  Siddharth Chandrasekar S (3122235001310)\\
  Vedanth M S (3122235001314)}
\supervisor{Dr. Anusha Jayasimhan\\Assistant Professor}
\department{Department of Computer Science and Engineering}
\institution{Sri Sivasubramaniya Nadar College of Engineering}
\academicyear{2026--2027}
\collegelogo{ssn-logo.png}

\begin{document}
\makereviewtitle
\reviewfrontmatter

\begin{reviewabstract}
ECG foundation models are increasingly used for automated interpretation and clinical
decision support, making their reliability under incomplete physiological observations an
important concern. In real-world acquisition and monitoring, ECG leads may be missing,
corrupted, or unavailable, yet a model may still produce confident interpretations from
insufficient evidence. Existing foundation and knowledge-enhanced approaches perform
strongly in representation learning, diagnosis, partial-lead analysis, and structured
interpretation, but are evaluated mainly through diagnostic accuracy and representation
quality. Consequently, how evidence-grounded reasoning degrades as physiological
information is removed remains insufficiently characterised.

This project investigates a knowledge-aware ECG reasoning framework under progressive
physiological information loss. Raw twelve-lead recordings are masked at seven availability
levels---12, 10, 8, 6, 4, 2 and 1---without reconstructing missing signals. Each degraded
recording is encoded by a frozen ECG foundation model into a latent representation,
converted into structured clinical findings, and passed to an evidence-based reasoning layer.
Matched conditions compare ECG-only reasoning with reasoning augmented by a Clinical
Knowledge Prior derived from SNOMED CT and ECG-specific clinical relationships.
The resulting diagnoses, supporting and opposing evidence, confidence, and uncertainty
are independently verified against an evidence schema and the lead-availability mask.
PTB-XL supports framework development, MIMIC-IV-ECG large-scale validation, and
Chapman-Shaoxing external validation.

Evaluation will measure grounding quality, evidence observability, diagnostic accuracy,
calibration, uncertainty, the Area Under the Reasoning Degradation Curve, and the
Reasoning Failure Boundary. As experiments are not complete, quantitative results
are not reported. The main contribution is a reproducible framework for locating the
boundary at which an ECG foundation model can no longer maintain evidence-grounded
reasoning and for testing whether structured clinical knowledge shifts that boundary
without promoting unsupported inference.

\keywords{Electrocardiogram; Foundation models; Evidence grounding; Reasoning
failure boundary; Clinical knowledge prior}
\end{reviewabstract}

\chapter{Introduction}

\subsection*{Project Domain}

This project lies at the intersection of cardiovascular informatics, deep learning
foundation models, and machine-learning safety in healthcare. Its primary focus is the
automated interpretation of multi-lead electrocardiograms (ECGs) and, more importantly,
the verification of whether an artificial intelligence system's clinical reasoning remains
grounded in the physiological evidence that is actually available to it. The work therefore
extends beyond conventional diagnostic classification: it studies the reliability,
traceability, and safety of model-generated clinical claims under progressively degraded
and resource-constrained acquisition conditions.

The proposed framework treats ECG interpretation as an evidence-dependent reasoning
task. A diagnostic conclusion is considered trustworthy only when the observations used
to support it can be verified from the available leads. This perspective is particularly
important in healthcare, where a plausible prediction is not sufficient unless the system
can demonstrate that its conclusion is consistent with the signals observed from the
patient. By combining physiological signal processing, foundation-model representations,
structured clinical observations, and deterministic verification, the project aims to
establish measurable safety limits for ECG-based clinical artificial intelligence.

\subsection*{Context}

The standard 12-lead ECG is one of the most widely used non-invasive tools for assessing
the electrical activity of the heart. Although it is referred to as a 12-lead system, it is
typically acquired using ten electrodes that produce twelve distinct electrical views.
The limb leads observe cardiac electrical activity primarily in the frontal plane, whereas
the precordial leads provide complementary views across the horizontal plane. Together,
these leads allow clinicians to examine cardiac rhythm, conduction abnormalities,
chamber-related changes, and regional patterns associated with conditions such as
myocardial ischaemia or infarction.

Recent advances in deep learning have enabled large-scale ECG foundation models to learn
general-purpose representations from extensive collections of labelled and unlabelled
recordings. Models such as ECG-FM, K-MERL, and GEM demonstrate strong performance on
tasks including rhythm classification, abnormality detection, and clinical report
generation when complete, high-quality ECG recordings are available. These developments
indicate the potential of foundation models to support clinical decision-making across a
wide range of cardiovascular applications.

In practice, however, a complete and artifact-free 12-lead recording cannot always be
obtained. Emergency departments, pre-hospital transport services, intensive-care units,
ambulances, and wearable monitoring systems frequently operate under acquisition
constraints. Electrodes may become detached, signals may be corrupted by motion, and
time-critical triage may limit the number of leads that can be collected. Consumer and
portable devices may also use reduced configurations consisting of six, four, two, or even
a single lead. These conditions create \emph{progressive physiological information loss}:
as leads are removed, the model loses not only signal samples but also specific anatomical
views of cardiac activity.

The effect of lead loss is therefore not equivalent to ordinary random noise or missing
tabular data. Each lead contributes a distinct spatial projection, and the removal of a
lead may make particular clinical findings partially observable or entirely unavailable.
A safe ECG reasoning system must recognise this change in observability and adjust its
claims, confidence, and uncertainty accordingly.

\subsection*{Motivation}

Most existing ECG artificial-intelligence benchmarks evaluate models using final-label
metrics such as accuracy, sensitivity, specificity, F1-score, or area under the receiver
operating characteristic curve. These measures are useful for assessing predictive
performance, but they do not determine whether the reasoning used to reach a prediction
is physically supported by the active sensor configuration. Evaluation is also commonly
performed under a fixed lead setting, which provides limited insight into how a model's
reliability changes as physiological information is progressively removed.

This limitation becomes critical when the evidence required for a diagnosis is regional.
For example, confirmation of an inferior myocardial infarction conventionally depends on
observing compatible ST-segment changes across contiguous inferior leads, particularly
Leads II, III, and aVF. If key inferior leads such as III and aVF are unavailable, the
regional evidence needed to confirm the finding is no longer fully observable. A model may
nevertheless continue to output the correct diagnostic label because of correlations learned
from the training data, patient-level patterns, or dataset-specific shortcuts. It may also
generate a clinical explanation that refers to changes in leads that were not provided as
input.

Conventional label-based evaluation can mistakenly reward such an output because the
predicted diagnosis matches the reference label. From a clinical-safety perspective,
however, a correct conclusion supported by unavailable or fabricated evidence is not
equivalent to a properly grounded conclusion. This behaviour creates a risk of silent
failure: the output appears plausible and confident even though the model has crossed the
limit at which the available physiological observations can justify its reasoning. In a
clinical environment, such unsupported claims may mislead practitioners, conceal the need
for additional measurements, and reduce trust in decision-support systems.

The project is motivated by the need to distinguish diagnostic correctness from reasoning
validity. Rather than asking only whether a model predicts the expected label, the proposed
work asks whether every supporting or opposing clinical claim is observable, consistent,
and verifiable under the current lead configuration. It further investigates whether a
structured Clinical Knowledge Prior (CKP) can improve reasoning under limited information
without encouraging the model to replace missing physiological evidence with unsupported
knowledge-based assumptions.

\subsection*{Significance}

This project addresses an important gap in trustworthy medical artificial intelligence by
formalising the point at which ECG reasoning changes from evidence-grounded deduction to
unsupported inference. The proposed \emph{Reasoning Failure Boundary} (RFB) identifies
the level of lead availability beyond which the system can no longer satisfy predefined
grounding and reliability requirements. In addition, the \emph{Area Under the Reasoning
Degradation Curve} (AURDC) summarises how reasoning quality changes across the complete
sequence of lead-degradation conditions, rather than reporting performance at only one
operating point.

The framework also introduces an explicit three-state verification mechanism that labels
each clinical claim as \textsc{supported}, \textsc{contradicted}, or \textsc{unavailable}.
This distinction prevents the absence of evidence from being incorrectly interpreted as
negative evidence. It also enables the system to state when a claim cannot be assessed from
the remaining leads, which is a safer and more informative response than producing an
unverifiable explanation.

By incorporating structured relationships derived from SNOMED CT through the CKP, the
study evaluates whether clinical knowledge can help preserve coherent reasoning as sensor
information decreases. Crucially, the knowledge prior is intended to organise and constrain
reasoning, not to reconstruct missing ECG signals or invent observations. The resulting
framework can provide clinicians and system developers with interpretable evidence trails,
calibrated uncertainty, and empirically measured operating limits. Such capabilities are
relevant to hospital-based ECG systems, emergency and pre-hospital care, remote monitoring,
and wearable devices, where the number and quality of available leads may vary over time.

More broadly, the work contributes a reproducible methodology for evaluating foundation
models in safety-critical settings where information is progressively lost. Its emphasis on
observability, claim-level verification, and explicit failure boundaries supports the
development of clinical AI systems that recognise when they possess sufficient evidence to
reason and when they must defer, request additional measurements, or recommend expert
review.

\section{Objectives}
\begin{enumerate}[label=\arabic*.,leftmargin=*,align=left]
  \item To design and develop a controlled degradation framework and an architecture
  with two complementary representations that model progressive physiological information loss
  across standard ECG configurations containing 12, 10, 8, 6, 4, 2, and 1 available
  leads. The framework will remove leads without introducing synthetic or reconstructed
  signals and will combine continuous latent feature embeddings with discrete,
  deterministic clinical observations.

  \item To implement and evaluate a structured reasoning pipeline with a deterministic
  three-state grounding-verification mechanism that classifies individual clinical claims
  as \textsc{supported}, \textsc{contradicted}, or \textsc{unavailable}. This mechanism
  will be used to determine whether the model's findings, evidence, and conclusions remain
  truthfully anchored in the physical sensor data at every degradation level.

  \item To compare an unaugmented ECG-only baseline with a knowledge-enhanced model that
  integrates a Clinical Knowledge Prior derived from the SNOMED CT clinical ontology into
  a shared reasoning transformer. The comparison will quantify changes in the Area Under
  the Reasoning Degradation Curve and determine whether the knowledge prior delays the
  Reasoning Failure Boundary under severe sensor scarcity without promoting unsupported
  clinical inference.
\end{enumerate}


\chapter{First Review: Queries and Responses}

The first project review raised important questions concerning the clinical validity of the
verification process, the necessity of the proposed dual-condition architecture, and the
practical value of the Reasoning Failure Boundary. This chapter records the questions and
provides the corresponding technical responses. The feedback has been incorporated into
the methodology by strengthening clinical oversight of the Evidence Schema, clarifying the
role of the control and experimental conditions, and defining how the proposed safety
measures can support real-world decision-making.

\section{Medical Reliability of the Evidence Schema}

\textbf{Query 1: How will you ensure that the Evidence Schema in Module 5 is medically
reliable?}

\textbf{Response:} The Evidence Schema will be designed and validated with guidance from
a qualified medical professional who possesses appropriate knowledge of ECG
interpretation. Clinical involvement is essential because Module 5 does not merely compare
text labels; it determines whether a model-generated claim is supported by the
physiological evidence available in a particular lead configuration. An incorrectly defined
relationship between a clinical finding and its required leads could therefore produce a
misleading grounding verdict even if the software implementation were technically correct.

For every admissible ECG claim, the schema will specify the clinical finding being assessed,
the waveform measurement required to assess it, the anatomical lead or lead group from
which that measurement must be obtained, and the decision criterion used to classify the
claim. The medical expert will assist in identifying clinically relevant findings, confirming
their relationships to cardiac diagnoses, and establishing the conditions under which each
finding may be considered supported, contradicted, partially observable, or unavailable.
Particular attention will be given to regional findings, such as ST-segment changes, whose
interpretation depends on the presence of specific contiguous leads.

The schema-development process will include the following stages:

\begin{enumerate}[label=\arabic*.,leftmargin=*,align=left]
  \item Construct an initial machine-readable schema using established ECG terminology,
  SNOMED CT concepts, and accepted clinical relationships.

  \item Review each claim type, required lead set, measurement definition, and verification
  criterion with the medical expert.

  \item Test the schema on representative complete-lead and reduced-lead ECG examples to
  confirm that supported, contradicted, and unavailable verdicts are assigned correctly.

  \item Record all approved schema elements, revisions, and clinical justifications under
  version control so that the verification procedure remains auditable and reproducible.
\end{enumerate}

The Evidence Schema will be fixed before final model evaluation. This prevents the
verification rules from being modified after observing experimental results and ensures
that the grounding assessment remains independent of the reasoning model. Through
clinical review, explicit lead dependencies, representative-case testing, and controlled
versioning, the schema will be medically meaningful as well as computationally
reproducible.

\section{Need for the Control and Experimental Conditions}

\textbf{Query 2: Why is the model complex? Do we actually need both the control and
experimental parts?}

\textbf{Response:} Both reasoning conditions are not required to construct the core
reasoning-degradation benchmark. The ECG-only control path is the essential component
because it establishes how the foundation model's reasoning changes as physiological
information is removed. It produces the baseline reasoning degradation curve from which
Grounding Quality, the Area Under the Reasoning Degradation Curve, and the Reasoning
Failure Boundary are calculated.

The experimental CKP path addresses a separate research question: whether structured
clinical knowledge can help preserve valid reasoning under reduced-lead conditions without
encouraging unsupported inference. The Clinical Knowledge Prior does not reconstruct
missing leads or supply artificial waveform observations. Instead, it provides structured
relationships among diagnoses, clinical findings, and anatomical dependencies that may
help the reasoning component organise the evidence still available to it.

Although the complete architecture contains several modules, its complexity serves
experimental control and auditability. Signal preprocessing, controlled masking, latent
representation, observation extraction, reasoning, and grounding verification are separated
so that failures can be localised to a specific stage. The two reasoning conditions share
the same patient recordings, masks, frozen ECG encoder, latent representations, structured
observations, candidate diagnoses, transformer architecture, and verification rules. They
differ only through the inclusion of the CKP in the experimental context. This matched
design makes the comparison interpretable: any systematic change in grounding performance
can be associated primarily with the addition of structured clinical knowledge rather than
with a change in input data or model capacity.

The control path will therefore remain the required benchmark, while the CKP path will be
treated as an experimental extension. Results from the control path are sufficient to locate
the baseline safety boundary. Comparison with the experimental path determines whether
the CKP improves AURDC or shifts the RFB toward more severely degraded lead
configurations. Presenting both results also prevents an improvement from being attributed
to the knowledge prior without a directly comparable baseline.

\section{Purpose and Clinical Utility of the Reasoning Failure Boundary}

\textbf{Query 3: What are you achieving with this work? How will the Reasoning Failure
Boundary be useful?}

\textbf{Response:} The project aims to determine how much observable ECG information is
required for a clinical AI system to produce reasoning that remains supported by the
physical sensor data. Existing evaluation commonly reports whether a predicted diagnosis
matches the reference label. This project additionally evaluates whether the clinical claims
used to justify that diagnosis can be verified from the leads that were actually available.
It therefore distinguishes diagnostic correctness from evidence-grounded reliability.

The Reasoning Failure Boundary provides an empirical operating limit. For a predefined
clinical-safety threshold, it identifies the most severely degraded tested lead configuration
at which the system's Grounding Quality remains acceptable. For example, if the model
meets the grounding threshold with six available leads but fails to meet it at four leads,
the boundary indicates that the model should not be assumed to provide sufficiently
evidence-grounded reasoning below the six-lead condition. The complete reasoning
degradation curve and AURDC are reported alongside the boundary so that this conclusion
does not depend on a single operating point.

This boundary can assist clinicians and clinical-AI developers in several ways. It can
indicate when additional electrodes or a complete 12-lead acquisition should be requested,
when the system should reduce its confidence or abstain from a diagnosis, and when the
output requires review by a qualified medical professional. It can also support the
selection of safe operating configurations for ambulances, emergency triage systems,
intensive-care monitoring, and wearable devices that do not always provide all twelve
leads.

The RFB is not intended to replace clinical judgement or certify a model for unrestricted
deployment. Its value is that it converts an otherwise hidden failure mode into a measurable
and interpretable safety indicator. By comparing the RFB of the control and CKP-assisted
conditions, the project can further determine whether structured knowledge extends the
range over which reliable reasoning is possible, while the claim-level verifier ensures that
such an extension is not achieved through hallucinated evidence.


\chapter{Problem statement}

Modern ECG foundation models can maintain apparently strong diagnostic performance even
when a substantial portion of the original sensor information is unavailable. Under
progressive lead loss, a model may continue to predict the correct diagnostic label by
exploiting statistical regularities, dataset-specific shortcuts, or correlations that do not
constitute direct physiological evidence for the prediction. More critically, a generative
reasoning component may produce clinical claims referring to waveform changes in leads
that were not supplied to the model. Conventional classification metrics do not distinguish
between a diagnosis supported by observable electrodynamic evidence and a plausible guess
supported by an unavailable or hallucinated finding.

This limitation represents a significant safety concern in emergency triage, pre-hospital
transport, intensive-care monitoring, and wearable telemetry, where missing or detached
electrodes are common. Many cardiac conditions are identified using patterns distributed
across anatomically related leads. When the required projections are absent, the available
recording may no longer contain sufficient information to confirm or reject a localised
pathology. A reliable clinical system must therefore recognise not only what diagnosis is
probable, but also whether the evidence needed to justify that diagnosis can be observed
from the active physical sensors.

The problem addressed in this project is the absence of a systematic method for measuring
how evidence-grounded clinical reasoning degrades as ECG leads are progressively removed.
The project seeks to determine the operating point at which a foundation model's reasoning
can no longer be considered sufficiently grounded, even if its final diagnostic accuracy
remains high. It also investigates whether structured clinical knowledge can delay this
failure point without replacing missing signals with artificial evidence or encouraging
unsupported inference.

Formally, the study considers the ordered set of lead-availability levels

\begin{equation}
  \mathcal{L} = (12, 10, 8, 6, 4, 2, 1),
\end{equation}

where each value denotes the number of physical ECG leads retained at a particular
degradation stage. The same patient recording is evaluated across these controlled stages,
allowing diagnostic performance, evidence observability, grounding quality, and uncertainty
to be measured as functions of the available physiological information. Two matched
reasoning conditions are compared: an ECG-only baseline and a knowledge-guided condition
augmented with a Clinical Knowledge Prior (CKP) derived from SNOMED CT.

\section{Inputs}

The proposed system receives four principal categories of input.

\subsection{Multi-lead ECG waveform}

The primary input is a multi-channel surface ECG recording represented as a voltage-time
matrix. Every recording is standardised to a fixed duration and sampling frequency so that
all degradation conditions are evaluated over the same temporal interval. The complete
recording contains the conventional twelve ECG leads. Signal values are retained only for
the leads that are physically available at the selected degradation level; unavailable leads
are not regenerated, inferred, or replaced with synthetic waveforms.

Before model inference, the retained signals undergo a consistent preprocessing pipeline
for quality control, amplitude normalisation, baseline-wander removal, and suppression of
high-frequency interference. The preprocessing operations must preserve diagnostically
important waveform morphology and temporal relationships, including the P wave, QRS
complex, ST segment, and T wave. Every model condition is supplied with the same processed
recording for a given patient and degradation level to ensure a controlled comparison.

\subsection{Lead-availability mask}

Each ECG recording is accompanied by a binary lead-availability mask

\begin{equation}
  \mathbf{m}^{(\ell)} = (m_1,m_2,\ldots,m_{12}), \qquad
  m_j \in \{0,1\}, \quad \ell \in \mathcal{L},
\end{equation}

where $m_j=1$ indicates that lead $j$ is physically available and $m_j=0$
indicates that it has been removed at the selected degradation level. The mask is treated as
an explicit representation of sensor observability rather than as a request to impute the
missing channel. It is propagated through the reasoning and verification stages so that
every clinical claim can be checked against the anatomical views present in the input.

The degradation schedule is deterministic and reproducible. Lead subsets are selected to
represent clinically meaningful reduced-lead configurations wherever possible, rather than
arbitrary combinations that have no practical acquisition interpretation. This enables the
resulting safety measurements to be related to realistic hospital, ambulatory, and wearable
monitoring settings.

\subsection{Candidate diagnostic hypotheses}

The system receives a predefined candidate set covering major cardiac disease and rhythm
categories represented in the evaluation datasets. These candidates define the diagnostic
space over which posterior probabilities are estimated. The same candidate set is used for
the baseline and knowledge-guided conditions, preventing changes in the output vocabulary
from influencing the comparison.

Each diagnostic hypothesis is associated with clinical findings that may support, oppose,
or be insufficient to assess the diagnosis. Where a finding depends on a particular
anatomical lead group, this dependency is represented explicitly so that the system can
distinguish between a negative finding and a finding that cannot be evaluated because the
required leads are unavailable.

\subsection{Clinical Knowledge Prior}

The knowledge-guided condition additionally receives a structured CKP derived from
SNOMED CT and ECG-specific clinical relationships. The CKP represents associations among
diagnoses, observable ECG findings, anatomical regions, and supporting or opposing evidence.
Its role is to constrain and organise the reasoning process by supplying clinically valid
relationships; it is not used to synthesise missing signals or declare an observation that is
not present in the recording.

For experimental control, both the baseline and CKP-enhanced systems use the same ECG
representation, diagnostic candidates, degradation masks, reasoning architecture, and data
splits. The presence or absence of the knowledge prior is the principal variable separating
the two conditions.

\section{Outputs}

For every recording and lead-availability level, the system produces diagnostic,
interpretability, verification, and degradation outputs.

\subsection{Diagnostic posterior distribution}

The model outputs a ranked posterior probability distribution over the candidate cardiac
diagnoses. In addition to the highest-ranked diagnosis, the complete distribution is retained
to assess diagnostic discrimination, calibration, and changes in confidence as leads are
removed. A safe response may include increased uncertainty or diagnostic abstention when
the remaining evidence is insufficient.

\subsection{Structured clinical reasoning trace}

The model generates an interpretable reasoning trace composed of discrete clinical claims.
Each claim identifies a physiological observation, its relationship to one or more candidate
diagnoses, whether it functions as supporting or opposing evidence, and an associated
evidence-attention weight. The trace is structured so that claims can be evaluated
individually rather than judged only through the fluency of a free-text explanation.

This representation makes it possible to detect cases in which the final label is correct but
the explanation is unsupported. It also permits direct comparison of the evidence selected
by the ECG-only and CKP-enhanced reasoning conditions at every degradation stage.

\subsection{Grounding-verification verdicts}

A deterministic verifier assigns every generated claim to exactly one of three states:

\begin{itemize}
  \item \textsc{Supported}: the required leads are available and the claimed physiological
  pattern is present in the observed signal;

  \item \textsc{Contradicted}: the required leads are available, but the observed waveform
  conflicts with the generated claim; or

  \item \textsc{Unavailable}: the claim cannot be evaluated because one or more of the
  anatomical projections required to verify it are absent.
\end{itemize}

The three-state design is essential because unavailable evidence is neither positive nor
negative evidence. A claim classified as \textsc{Unavailable} reveals that the system has
attempted to reason beyond the observable sensor space and must therefore express
uncertainty rather than present the claim as an established finding.

\subsection{Reasoning-degradation metrics}

The claim-level verdicts are aggregated at each lead level to produce the following
quantitative outputs:

\begin{itemize}
  \item \textbf{Evidence Observability}, which measures the proportion of diagnostic
  evidence that can be assessed using the currently available lead configuration;

  \item \textbf{Grounding Quality}, which measures the extent to which the generated
  reasoning is consistent with observable physiological evidence while penalising
  contradicted and unjustifiably asserted unavailable claims;

  \item \textbf{True Grounding Ratio}, which measures the proportion of generated claims
  that are explicitly verified as supported by the retained ECG signals;

  \item \textbf{Area Under the Reasoning Degradation Curve (AURDC)}, a normalised summary
  of Grounding Quality across all seven ordered lead-availability levels; and

  \item \textbf{Reasoning Failure Boundary (RFB)}, the first degradation stage at which
  Grounding Quality falls below the specified clinical-safety threshold.
\end{itemize}

These outputs are reported together with conventional diagnostic and calibration metrics.
This prevents a high classification score from concealing a decline in evidence-grounded
reasoning.

\section{Constraints}

The investigation is conducted under the following operational and scientific constraints.

\subsection{No reconstruction of unavailable leads}

Generative lead reconstruction, deep inpainting, mathematical lead synthesis, and
substitution using predicted waveforms are prohibited. A missing lead must remain missing
throughout the experiment. This restriction ensures that the system is evaluated under
genuine physiological information scarcity and prevents artificially generated signal
content from being treated as observed clinical evidence.

\subsection{Morphology-preserving signal preprocessing}

Signal filtering must reduce baseline wander and high-frequency noise without shifting or
distorting diagnostically relevant waveform intervals. Accordingly, digital filters are
applied in a zero-phase, bidirectional manner. Forward--backward filtering cancels net
phase delay and prevents preprocessing-induced displacement of waveform features such as
QRS boundaries or ST-segment changes. Identical preprocessing parameters are used across
all model conditions and lead levels.

\subsection{Patient-stratified evaluation}

Training, validation, and testing partitions must be created at the patient level. All ECG
recordings belonging to the same patient are restricted to a single partition. This prevents
patient-overlap leakage, in which subject-specific morphology observed during training
could reappear during evaluation and produce an unrealistically favourable estimate of
generalisation.

\subsection{Controlled experimental comparison}

The ECG-only baseline and CKP-enhanced condition must be evaluated using identical
recordings, lead masks, candidate diagnoses, preprocessing operations, and verification
rules. Lead removal is performed according to the same ordered schedule for both systems.
No condition may receive patient information or signal channels that are unavailable to the
other condition. Any observed difference can therefore be attributed primarily to the
inclusion of structured clinical knowledge.

\subsection{Evidence-limited use of clinical knowledge}

The CKP may represent clinical relationships and dependencies, but it must not override the
lead-availability mask or convert an unobservable finding into a supported claim. Knowledge
can guide the selection and organisation of evidence; it cannot substitute for physical
measurement. This constraint is necessary to determine whether the knowledge prior
improves grounded reasoning rather than merely producing more persuasive explanations.

\section{Success Criteria}

The proposed system is evaluated using explicit and measurable safety criteria. Let
$GQ(\ell)$ denote the Grounding Quality measured with $\ell$ available leads, and
let the clinical-safety threshold be

\begin{equation}
  \tau_{\mathrm{safety}} = 0.60.
\end{equation}

For the tested lead levels $\mathcal{L}=\{12,10,8,6,4,2,1\}$, the RFB is defined as
the lowest lead availability at which the grounding score still satisfies the threshold:

\begin{equation}
  \mathrm{RFB}=\min\{\ell\in\mathcal{L}:GQ(\ell)\geq 0.60\}.
\end{equation}

The resulting value identifies the most severely degraded tested configuration for which
the model's reasoning remains clinically acceptable. A lower RFB therefore indicates
greater tolerance to physiological information loss. If no evaluated lead level satisfies
the threshold, the RFB is reported as undefined.

The knowledge-guided system will be considered successful when it satisfies both of the
following primary criteria:

\begin{enumerate}[label=\arabic*.,leftmargin=*,align=left]
  \item It achieves a normalised AURDC of at least $0.80$ across the full sequence from
  twelve available leads to a single available lead, demonstrating sustained reasoning
  fidelity rather than strong performance at only one lead configuration.

  \item Its RFB reduces the minimum acceptable lead requirement by at least two physical
  leads relative to the unaugmented ECG-only baseline. In practical terms, the CKP-enhanced
  model must remain at or above the $0.60$ Grounding Quality threshold at a configuration
  containing at least two fewer active leads than the lowest acceptable baseline
  configuration.
\end{enumerate}

Secondary evaluation examines whether this improvement is achieved without a clinically
unacceptable reduction in diagnostic discrimination or calibration. Results are also
reported separately for the \textsc{Supported}, \textsc{Contradicted}, and
\textsc{Unavailable} claim states so that an increased AURDC cannot be attributed to
unrestricted use of prior knowledge. Together, these criteria define success as the
preservation of verifiable, evidence-grounded reasoning under sensor scarcity, rather than
the preservation of diagnostic label accuracy alone.

\chapter{Methodology}

The proposed framework is organised into six sequential phases. An ECG recording is
first ingested and standardised, progressively degraded using controlled lead masking,
encoded into a latent representation by a frozen foundation model, processed under two
parallel clinical-reasoning conditions, verified claim by claim against the physiological
evidence that remains observable, and finally summarised using the project’s degradation
and failure-boundary measures. Each phase consumes a defined artefact from its predecessor
and produces an auditable output for the next phase. This modular organisation allows the
pipeline to be implemented, validated, and reproduced in dependency order.

\section{Architectural Design for Proposed System}

Figure~\ref{fig:architecture} presents the high-level architecture of the proposed system.
The design contains six connected phases: data preparation, controlled masking, latent
representation, clinical reasoning, grounding verification, and results generation. The
architecture also exposes a direct connection from the masking phase to the grounding
phase. This path is essential because the verifier must receive the true lead-availability
mask and the masked signal directly, rather than relying on the reasoning model’s own
description of what it observed.

\begin{figure}[H]
  \centering
  \includegraphics[width=\textwidth]{../Gemini_Generated_Image_e5sdgje5sdgje5sd.png}
  \caption{High-level architecture of the proposed system}
  \label{fig:architecture}
\end{figure}

The first phase covers data and ground truth. Recordings from PTB-XL,
MIMIC-IV-ECG, and Chapman-Shaoxing pass through five ordered preprocessing
operations: canonical lead ordering, sampling-rate and duration standardisation, unit
calibration and normalisation, filtering, and quality screening. The phase emits a
standardised twelve-lead tensor. Diagnostic annotations are separated from the signal and
retained exclusively as ground truth so that they cannot leak into the model input.

The second phase applies controlled physiological degradation. Progressive lead masking
produces seven variants containing 12, 10, 8, 6, 4, 2, and 1 available leads. Every variant
is emitted with a binary availability mask identifying the retained physical channels.
The operation removes leads without reconstructing or imputing their contents.

The third phase generates the signal representation. Each masked recording and its
availability mask are presented to ECG-FM. The foundation-model weights remain frozen
throughout the experiment, and the resulting latent representation is cached at every lead
level. Both downstream reasoning conditions therefore receive the same encoded signal,
ensuring that differences between them cannot originate from the encoder.

The fourth phase performs clinical reasoning under two matched conditions. A structured
observation extractor converts the latent representation into ECG findings covering
rhythm, rate, PR interval, QRS duration, and ST-segment changes. A shared rule base uses
these observations to construct the same candidate hypothesis set for both conditions.
The control condition reasons from the observations and candidate evidence alone. The
experimental condition additionally receives a Clinical Knowledge Prior (CKP) derived
from a curated SNOMED CT knowledge graph. A shared transformer produces candidate
scores and evidence attribution, yielding a structured reasoning trace.

The fifth phase verifies grounding. The reasoning trace is decomposed into individually
adjudicable clinical claims. Each claim is compared with an independently defined
evidence schema, the masked recording, and the availability mask. The verifier determines
whether the claim is supported, contradicted, or unavailable, and the results are
aggregated into Evidence Observability and Grounding Quality.

The sixth phase produces the reported results. Grounding Quality across all seven lead
levels forms the reasoning degradation curve. The Area Under the Reasoning Degradation
Curve (AURDC) and the Reasoning Failure Boundary (RFB) are computed separately for
the control and CKP-assisted conditions. Their paired difference constitutes the primary
experimental comparison. Any numerical values depicted in the results block of
Figure~\ref{fig:architecture} are illustrative elements of the architecture and are not
experimental findings.

\subsection*{End-to-end evaluation procedure}

For every ECG recording, the complete evaluation is performed as follows:

\begin{enumerate}[label=\textbf{Step \arabic*:},leftmargin=*,align=left]
  \item Standardise the raw recording using the ordered preprocessing pipeline and retain
  the diagnostic annotation separately as ground truth.

  \item For each lead level in the ordered set
  $\Lambda=(12,10,8,6,4,2,1)$, generate a masked recording and its corresponding
  binary lead-availability mask.

  \item Encode the masked recording using the frozen ECG foundation model and cache the
  latent representation.

  \item Extract structured ECG observations and construct a candidate diagnostic
  hypothesis set that is shared by both reasoning conditions.

  \item Generate one reasoning trace using the ECG-only control context and a second trace
  using the CKP-assisted context.

  \item Extract the claims in each trace and verify every claim against the evidence schema,
  masked waveform, and lead-availability mask.

  \item Aggregate the claim verdicts into Evidence Observability and Grounding Quality for
  each lead level and reasoning condition.

  \item Construct the reasoning degradation curves and compute AURDC, RFB, and the paired
  improvement produced by the CKP-assisted condition.
\end{enumerate}

\section{Modules}

The framework consists of six modules. Each module is defined by its purpose, inputs,
processing operations, outputs, and dependencies. Let

\begin{equation}
  \mathcal{C}=\{\mathrm{I},\mathrm{II},\mathrm{III},\mathrm{aVR},
  \mathrm{aVL},\mathrm{aVF},\mathrm{V1},\ldots,\mathrm{V6}\}
\end{equation}

denote the standard ECG lead set, and let
$\Lambda=(12,10,8,6,4,2,1)$ denote the evaluated lead-availability levels. A
standardised recording is represented by
$\mathbf{X}\in\mathbb{R}^{12\times n_T}$, where $n_T$ is the number of temporal
samples. At level $l\in\Lambda$, $\mathbf{a}_l\in\{0,1\}^{12}$ is the availability
mask and

\begin{equation}
  \mathcal{A}_l=\{\lambda\in\mathcal{C}:\mathbf{a}_l[\lambda]=1\}
\end{equation}

is the set of leads physically available to the model and verifier.

\subsection{Module 1: Data and Input Management}

\textbf{Purpose.} To establish a reproducible data pathway and strictly separate the
physiological signal that the model is permitted to observe from the annotations used to
evaluate its output.

\textbf{Inputs.} Raw recordings from PTB-XL, MIMIC-IV-ECG, and
Chapman-Shaoxing, together with their diagnostic annotations. The corpora contain
heterogeneous sampling frequencies, recording durations, lead orderings, amplitude units,
and device conventions.

\textbf{Processing.} Five ordered preprocessing operations are applied. First, lead
ordering enforces the canonical twelve-lead sequence; otherwise, later components that
address leads positionally could silently associate a finding with the wrong anatomical
projection. Second, sampling-rate and duration standardisation produce a common temporal
resolution and window length. Third, unit calibration converts stored device values into
physical millivolts using the record-specific gain and offset, after which a fixed
normalisation is applied. Fourth, zero-phase filtering suppresses baseline wander,
high-frequency interference, and power-line noise while preserving waveform timing.
Finally, quality screening rejects tensors that fail shape validation, contain non-finite
values, or include flatlined channels, and records the reason for every exclusion.

The preprocessing pathway is expressed as

\begin{equation}
  \mathbf{X}
  =(\mathcal{P}_5\circ\mathcal{P}_4\circ\mathcal{P}_3
  \circ\mathcal{P}_2\circ\mathcal{P}_1)(\mathbf{X}_{\mathrm{raw}}),
  \label{eq:preprocessing}
\end{equation}

where the operators respectively perform lead ordering, sampling and duration
standardisation, unit calibration and normalisation, filtering, and quality screening.
Equation~\ref{eq:preprocessing} fixes their execution order, while the concrete parameters
of every operator are retained in the record manifest.

\textbf{Outputs.} A standardised tensor of dimension $12\times n_T$, a ground-truth
annotation set held outside the model input, and a manifest containing the record
identifier, patient split, preprocessing parameters, and any rejection reason.

\textbf{Dependencies.} This is the initial module. All later modules depend on its
standardised signal and manifest.

\subsection{Module 2: Controlled Degradation}

\textbf{Purpose.} To create a controlled and measurable loss of physiological
observability. This module defines the principal independent variable of the study.

\textbf{Inputs.} The standardised twelve-lead tensor produced by Module 1.

\textbf{Processing.} A binary mask selects the leads retained at each of the seven
availability levels. Channels not selected are filled with a fixed missing-channel
representation, while the availability mask is retained as a separate input. This prevents
the placeholder value from being interpreted as physiological signal. No missing lead is
reconstructed or imputed, because doing so would convert the experiment from an
investigation of information loss into an evaluation of reconstruction quality.

With $\nu$ denoting the fixed missing-channel representation, masking is defined by

\begin{equation}
  \mathbf{X}_l
  =\mathcal{M}(\mathbf{X},\mathbf{a}_l)
  =\operatorname{diag}(\mathbf{a}_l)\mathbf{X}
  +\nu(\mathbf{1}_{12}-\mathbf{a}_l)\mathbf{1}_{n_T}^{\top},
  \qquad |\mathcal{A}_l|=l.
  \label{eq:masking}
\end{equation}

The second term in Equation~\ref{eq:masking} fills only unavailable rows and has
dimension $12\times n_T$.

\textbf{Outputs.} Seven degraded variants for every record,

\begin{equation}
  \{(\mathbf{X}_l,\mathbf{a}_l):l\in\Lambda\},
  \label{eq:masked-set}
\end{equation}

each paired with its explicit lead-availability mask.

\textbf{Dependencies.} Module 2 consumes the standardised signal from Module 1. It
supplies Module 3 with model input and supplies Module 5 directly with the masked signal
and mask required for independent verification.

\subsection{Module 3: Latent Representation}

\textbf{Purpose.} To convert each degraded recording into a representation of the
physiological information that remains observable.

\textbf{Inputs.} A masked ECG tensor and its lead-availability mask.

\textbf{Processing.} The tensor--mask pair is presented to ECG-FM with the foundation
model weights frozen. Freezing is a methodological control: adapting the encoder
separately at each lead level could allow it to compensate for the imposed degradation,
thereby confounding representation adaptation with the reasoning degradation that the
study intends to measure. The mask is supplied explicitly so that unavailable channels
are not interpreted as valid observations.

The representation step is

\begin{equation}
  \mathbf{Z}_l=f_{\theta}(\mathbf{X}_l,\mathbf{a}_l)
  \in\mathbb{R}^{D},\qquad \theta\ \text{fixed}.
  \label{eq:representation}
\end{equation}

Every representation is cached using the record identifier, lead level, and encoder
checkpoint. The control and CKP-assisted reasoning conditions consequently consume the
same $\mathbf{Z}_l$. Before clinical reasoning is introduced, sensitivity to masking is
audited using cosine divergence from the complete recording:

\begin{equation}
  \Delta^{\mathrm{cos}}_l
  =1-\frac{\langle\mathbf{Z}_{12},\mathbf{Z}_l\rangle}
  {\lVert\mathbf{Z}_{12}\rVert_2\lVert\mathbf{Z}_l\rVert_2}.
  \label{eq:cosine}
\end{equation}

This is a diagnostic audit rather than a primary outcome measure. It is not assumed to be
monotonic because different lead combinations with similar channel counts need not
contain the same diagnostic information.

\textbf{Outputs.} One cached ECG latent representation per record and lead level.

\textbf{Dependencies.} Module 3 consumes Module 2 and supplies Module 4. The shared
cache enforces an identical representation for both reasoning conditions.

\subsection{Module 4: Clinical Reasoning}

\textbf{Purpose.} To convert the latent representation into an inspectable reasoning
trace under two conditions that differ only in whether structured clinical knowledge is
available.

\textbf{Inputs.} The latent representation from Module 3, the evidence schema, the
shared hypothesis-generation rule base, and, in the experimental condition, the
SNOMED CT knowledge representation $\mathbf{Z}_K$.

\textbf{Processing.} Four stages are performed: observation extraction, hypothesis
generation, candidate-evidence preparation, and transformer-based reasoning.

The observation extractor projects the latent vector into logits for a fixed field set

\begin{equation}
  \mathcal{F}=\{\text{rhythm},\text{rate},\text{PR},\text{QRS},\text{ST}\}.
\end{equation}

For every field $f\in\mathcal{F}$,

\begin{align}
  \mathbf{z}^{(f)}_l &=\mathbf{W}_f\mathbf{Z}_l+\mathbf{b}_f,
  \label{eq:obs-logits}\\
  \widehat{y}^{(f)}_l
  &=\mathop{\mathrm{arg\,max}}\,
  \mathrm{softmax}\!\left(\mathbf{z}^{(f)}_l\right).
  \label{eq:obs-class}
\end{align}

A schema-validation layer admits only fields and values belonging to the fixed clinical
vocabulary. Every retained observation also records the lead set from which it was
derived. This intermediate perception layer allows errors of observation to be separated
from errors of inference and creates explicit claims that can later be verified.

The same rule base $\mathcal{R}$ converts the structured observation set
$\mathcal{O}_l$ into a shared candidate diagnosis set:

\begin{equation}
  \mathcal{H}_l=\mathcal{R}(\mathcal{O}_l)
  =\{h_1,\ldots,h_K\},\qquad
  \mathcal{H}^{(0)}_l=\mathcal{H}^{(1)}_l=\mathcal{H}_l,
  \label{eq:hypotheses}
\end{equation}

where condition $c=0$ denotes the ECG-only control and $c=1$ denotes the
knowledge-assisted condition. Equality of the candidate sets is an experimental
invariant.

For candidate $h_i$, the relevant SNOMED CT subgraph is denoted by
$\mathcal{G}_{l,i}$ and encoded by $\phi$. The knowledge representation is projected
into the ECG latent dimension and fused additively:

\begin{align}
  \mathbf{Z}^{K}_{l,i}
  &=\mathbf{W}_K\phi(\mathcal{G}_{l,i})+\mathbf{b}_K
  \in\mathbb{R}^{D},
  \label{eq:knowledge-projection}\\
  \mathrm{CKP}_{l,i}
  &=\mathbf{Z}_l+\mathbf{Z}^{K}_{l,i}.
  \label{eq:ckp}
\end{align}

Fusion occurs after hypothesis generation so that the knowledge prior cannot alter the
candidate set being compared. If $\mathcal{E}_{l,i}$ denotes the evidence associated
with $h_i$, the matched reasoning contexts are

\begin{align}
  \mathcal{C}^{(0)}_{l,i}
  &=[\mathcal{O}_l;\mathcal{E}_{l,i};h_i],
  \label{eq:control-context}\\
  \mathcal{C}^{(1)}_{l,i}
  &=[\mathcal{O}_l;\mathrm{CKP}_{l,i};\mathcal{E}_{l,i};h_i].
  \label{eq:experimental-context}
\end{align}

Thus, the two contexts differ by exactly one knowledge-derived term. A shared transformer
processes each context, while a score head estimates candidate support and an attention
head attributes evidence:

\begin{equation}
  \mathbf{U}^{(c)}_{l,i}
  =\mathrm{Trf}_{\varphi}\!\left(\mathcal{C}^{(c)}_{l,i}\right),
  \qquad
  s^{(c)}_{l,i}
  =\mathbf{w}_s^{\top}\mathbf{U}^{(c)}_{l,i}+b_s.
  \label{eq:transformer-score}
\end{equation}

Temperature scaling is fitted on validation data independently at every lead level:

\begin{equation}
  p^{(c)}_{l,i}
  =\frac{\exp(s^{(c)}_{l,i}/\beta_l)}
  {\sum_{i'=1}^{K}\exp(s^{(c)}_{l,i'}/\beta_l)},
  \qquad
  \widehat{d}^{(c)}_l
  =\mathop{\mathrm{arg\,max}}_{h_i\in\mathcal{H}_l}p^{(c)}_{l,i}.
  \label{eq:calibration}
\end{equation}

Grounding outcomes are excluded from calibration so that confidence and grounding remain
independently measurable.

\textbf{Outputs.} A structured trace for every record, lead level, and reasoning condition.
The trace contains the primary diagnosis, alternative diagnoses, calibrated probabilities,
supporting and opposing evidence, attention weights, and the lead attribution of every
cited item. Lead attribution is mandatory because Module 5 cannot adjudicate a claim
whose required anatomical evidence is unspecified.

\textbf{Dependencies.} Module 4 consumes Module 3, the evidence schema, and the
knowledge representation. It supplies Module 5.

\subsection{Module 5: Grounding Verification}

\textbf{Purpose.} To determine independently whether the reasoning produced by
Module 4 is supported by evidence that remained physically available at the current
masking level.

\textbf{Inputs.} The structured reasoning trace from Module 4; the masked recording and
availability mask supplied directly by Module 2; and the independently defined evidence
schema.

\textbf{Processing.} Claim extraction decomposes the reasoning trace into individually
adjudicable claims. For each claim $\gamma$, the evidence schema provides

\begin{equation}
  \mathcal{S}(\gamma)
  =\bigl(\mathrm{meas}(\gamma),\mathcal{R}(\gamma),\pi(\gamma)\bigr),
  \label{eq:evidence-schema}
\end{equation}

where $\mathrm{meas}(\gamma)$ identifies the required measurement,
$\mathcal{R}(\gamma)\subseteq\mathcal{C}$ is the required lead set, and
$\pi(\gamma)$ is the predefined verification criterion. The schema is fixed before
evaluation and is independent of the observation extractor, preventing the verifier from
measuring only the reasoning model’s self-consistency.

The fraction of required evidence that remains observable is

\begin{equation}
  o(\gamma,l)
  =\frac{|\mathcal{R}(\gamma)\cap\mathcal{A}_l|}
  {|\mathcal{R}(\gamma)|}\in[0,1].
  \label{eq:claim-observability}
\end{equation}

If $o(\gamma,l)=0$, the claim is labelled \textsc{Unavailable}. If
$o(\gamma,l)>0$, the verification criterion is evaluated on the observable evidence and
the claim is labelled \textsc{Supported} or \textsc{Contradicted}. The fractional value
is retained so that partially observable claims are not mistaken for fully observable
claims. Results are additionally stratified into fully observable claims, for which
$o(\gamma,l)=1$, and partially observable claims, for which
$0<o(\gamma,l)<1$.

For the extracted claim set $\Gamma_l$, the total required and observable evidence counts
are

\begin{align}
  N^{\mathrm{req}}_l
  &=\sum_{\gamma\in\Gamma_l}|\mathcal{R}(\gamma)|,
  \label{eq:required-evidence}\\
  N^{\mathrm{obs}}_l
  &=\sum_{\gamma\in\Gamma_l}
  |\mathcal{R}(\gamma)\cap\mathcal{A}_l|.
  \label{eq:observable-evidence}
\end{align}

Let $N^{\mathrm{sup}}_l$ and $N^{\mathrm{con}}_l$ denote the numbers of supported
and contradicted claims. The primary grounding metrics are

\begin{align}
  \mathrm{GQ}_l
  &=\frac{N^{\mathrm{sup}}_l}
  {N^{\mathrm{sup}}_l+N^{\mathrm{con}}_l},
  \label{eq:gq}\\
  \mathrm{EO}_l
  &=\frac{N^{\mathrm{obs}}_l}{N^{\mathrm{req}}_l}.
  \label{eq:eo}
\end{align}

Grounding Quality is undefined when no claim is evaluable. Such cases are excluded from
the level aggregate, and their count is reported. The evaluable-claim count
$N^{\mathrm{sup}}_l+N^{\mathrm{con}}_l$ is always reported alongside both ratios so
that a model cannot appear robust merely by emitting very few claims.

\textbf{Outputs.} A verification verdict and observability value for every claim,
together with per-level Evidence Observability and Grounding Quality.

\textbf{Dependencies.} Module 5 consumes Modules 2 and 4 and supplies Module 6.

\subsection{Module 6: Results and Evaluation}

\textbf{Purpose.} To convert the per-level grounding measurements into the project’s
principal reporting quantities and compare the ECG-only and CKP-assisted conditions.

\textbf{Inputs.} Evidence Observability, Grounding Quality, diagnostic performance,
calibration, and evaluable-claim counts at all seven lead levels for both conditions.

\textbf{Processing.} Grounding Quality is plotted against the number of available leads
to form the reasoning degradation curve:

\begin{equation}
  \mathcal{D}(l)=\mathrm{GQ}_l,\qquad l\in\Lambda.
  \label{eq:degradation-curve}
\end{equation}

Dispersion across records is shown alongside the aggregate curve. No monotonicity is
assumed because removing a lead may occasionally increase apparent grounding by reducing
the number of adjudicable claims; this interaction is reported rather than smoothed away.

Let $(L_1,\ldots,L_7)=(12,10,8,6,4,2,1)$. Because the lead levels are unevenly
spaced, the normalised AURDC is calculated using the weighted trapezoidal rule:

\begin{equation}
  \mathrm{AURDC}
  =\frac{1}{11}\sum_{i=1}^{6}
  \frac{\mathrm{GQ}_{L_i}+\mathrm{GQ}_{L_{i+1}}}{2}
  (L_i-L_{i+1}).
  \label{eq:aurdc}
\end{equation}

The factor $11=12-1$ normalises the area so that
$\mathrm{AURDC}\in[0,1]$ when all level-wise grounding scores lie in $[0,1]$.
Higher values indicate better preservation of evidence-grounded reasoning across the
complete degradation trajectory. Equation~\ref{eq:aurdc} is applied both to the aggregate
curve and to per-record curves. Paired per-record comparisons use only recordings with
defined Grounding Quality at every level; no missing score is imputed.

For a reliability threshold $\tau$ fixed before evaluation, the RFB is defined as

\begin{equation}
  \mathrm{RFB}
  =\min\{l\in\Lambda:\mathrm{GQ}_l\geq\tau\}.
  \label{eq:rfb}
\end{equation}

The RFB is the lowest tested lead availability at which grounding remains acceptable;
therefore, a lower value indicates greater tolerance to physiological information loss. If
the satisfying set is empty, the RFB is reported as undefined rather than assigned an
artificial value. The complete curve and a threshold-sensitivity analysis
$\mathrm{RFB}(\tau)$ are reported because a single boundary value can conceal
non-monotonic behaviour.

The effect of the Clinical Knowledge Prior is summarised without reversing the direction
of merit:

\begin{align}
  \Delta\mathrm{GQ}_l
  &=\mathrm{GQ}^{\mathrm{CKP}}_l-\mathrm{GQ}^{\mathrm{Control}}_l,
  \nonumber\\
  \Delta\mathrm{AURDC}
  &=\mathrm{AURDC}^{\mathrm{CKP}}
  -\mathrm{AURDC}^{\mathrm{Control}},
  \nonumber\\
  \Delta\mathrm{RFB}_{\mathrm{imp}}
  &=\mathrm{RFB}^{\mathrm{Control}}-\mathrm{RFB}^{\mathrm{CKP}}.
  \label{eq:condition-differences}
\end{align}

Positive values in Equation~\ref{eq:condition-differences} consistently indicate an
improvement produced by the knowledge-assisted condition. Comparisons are paired at the
record level and are reported with effect sizes and uncertainty intervals in addition to
significance tests.

\textbf{Outputs.} The reasoning degradation curve, AURDC, and RFB for each reasoning
condition, together with their paired differences and supporting diagnostic, calibration,
observability, and claim-count statistics.

\textbf{Dependencies.} Module 6 consumes the verified outputs of Module 5 and is the
terminal component of the proposed methodology.

\chapter{Implementation}

\section{Module 1: Data Acquisition, Preprocessing, and Controlled Degradation}

Module 1 implements the foundational data-engineering and signal-processing subsystem of
the framework. Its responsibility is to ingest multi-channel ECG recordings, remove or
control acquisition artefacts through a deterministic sequence of operations, and produce
standardised multi-lead arrays for controlled masking and downstream inference. The module
is implemented in Python~3.11 using NumPy, SciPy, Pandas, and the Waveform Database
(WFDB) library. Its architecture contains three principal stages: patient-stratified dataset
partitioning, a five-operation preprocessing pipeline, and a lead-availability masking
engine.

Every transformation is recorded in a per-record manifest. The manifest stores the source
identifier, patient identifier, dataset split, original and target sampling frequencies,
windowing or padding applied, calibration parameters, filtering configuration, quality
verdict, active lead set, and output tensor location. This makes the exact input supplied to
the model reproducible at every degradation level.

\subsection{Dataset Ingestion and Patient-Stratified Partitioning}

The primary development benchmark is PTB-XL version~1.0.3, which contains 21,799
clinical 10-second, 12-lead ECG recordings from 18,869 unique patients. Diagnostic
statements follow the SCP-ECG standard and are mapped to five principal diagnostic
superclasses: Normal ECG (\textsc{NORM}), Myocardial Infarction (\textsc{MI}),
ST/T Changes (\textsc{STTC}), Conduction Disturbance (\textsc{CD}), and
Hypertrophy (\textsc{HYP}).

The official fold assignments are retained, and partitioning is performed at the patient
level. All recordings belonging to a patient remain within one split, preventing the model
from observing the same patient-specific cardiac geometry in both development and
evaluation data. Table~\ref{tab:ptbxl-splits} summarises the resulting partitions.

\begin{table}[H]
  \centering
  \caption{Patient-stratified PTB-XL partitions used by the implementation}
  \label{tab:ptbxl-splits}
  \begin{tabularx}{\textwidth}{@{}lccc>{\raggedright\arraybackslash}X@{}}
    \toprule
    \textbf{Partition} & \textbf{Folds} & \textbf{Records} &
    \textbf{Patients} & \textbf{Purpose} \\
    \midrule
    Training & 1--8 & 17,418 & 15,096 &
    Adapter optimisation and representation learning \\
    Validation & 9 & 2,183 & 1,893 &
    Hyperparameter selection, calibration, and early stopping \\
    Held-out test & 10 & 2,198 & 1,880 &
    Quarantined reasoning-degradation benchmark \\
    \bottomrule
  \end{tabularx}
\end{table}

The held-out test partition is not used to select preprocessing parameters, optimise the
adapter, tune thresholds, or fit confidence calibration. It is accessed only after the
development configuration has been fixed. The distinction between recordings and patients
is retained throughout reporting: Fold~10 contains 2,198 recordings belonging to 1,880
patients.

\subsection{Five Sequential Preprocessing Operations}

Raw clinical ECGs vary in channel sequence, sampling frequency, amplifier gain, baseline
offset, and noise content. The \texttt{ECGPreprocessingPipeline} applies five operations
in a fixed order so that these variations do not become confounding variables. Both a
physical-unit signal and a model-normalised signal are retained: the former is used for
clinical rule verification, while the latter is used for latent representation learning.

\subsubsection{Operation P1: Lead Ordering and Spatial Alignment}

The \texttt{LeadOrdering} component maps source channel labels to the canonical sequence

\begin{equation}
  (\mathrm{I},\mathrm{II},\mathrm{III},\mathrm{aVR},\mathrm{aVL},
  \mathrm{aVF},\mathrm{V1},\mathrm{V2},\mathrm{V3},\mathrm{V4},
  \mathrm{V5},\mathrm{V6}).
\end{equation}

The limb leads describe cardiac electrical activity in the frontal plane, whereas the
precordial leads provide complementary views in the transverse plane. Preserving their
identity and order is therefore a clinical requirement rather than a formatting preference.
After mapping, the array is transposed into a channel-first layout of shape $(12,T)$. A
strict \texttt{LeadOrderingError} is raised if a required lead is absent, duplicated, or
cannot be mapped unambiguously.

\subsubsection{Operation P2: Sampling-Rate and Duration Standardisation}

Measurements such as the PR interval, QRS duration, and QT interval require a constant
temporal resolution. The \texttt{SamplingStandardizer} converts source recordings sampled
at rates such as 100~Hz or 1000~Hz to a common rate of 500~Hz using Fourier-domain
interpolation implemented by \texttt{scipy.signal.resample}.

After resampling, every recording is standardised to 10.0 seconds. Longer recordings are
cropped symmetrically about their temporal midpoint. Shorter recordings are padded
symmetrically with zeros, and the padded samples are recorded in a temporal validity mask
so that they cannot be treated as measured physiology. The final signal has the fixed
shape

\begin{equation}
  12\times(500\ \mathrm{samples/s}\times10\ \mathrm{s})
  =12\times5000.
\end{equation}

\subsubsection{Operation P3: Unit Calibration and Physical-Voltage Normalisation}

Stored digital counts depend on device-specific analogue-to-digital conversion gains and
baseline offsets. The \texttt{UnitCalibrator} converts each channel into millivolts:

\begin{equation}
  V_{\mathrm{mV}}[n]
  =\frac{x_{\mathrm{raw}}[n]-b_{\mathrm{ADC}}}{g_{\mathrm{ADC}}},
  \label{eq:adc-calibration}
\end{equation}

where $x_{\mathrm{raw}}[n]$ is the recorded count, $b_{\mathrm{ADC}}$ is the device
baseline, and $g_{\mathrm{ADC}}$ is the gain in counts per millivolt. Values outside the
physiologically plausible acquisition range $[-15,+15]$~mV are clipped and flagged for
the subsequent saturation audit.

Two synchronised outputs are then created. The calibrated millivolt array is preserved for
threshold-based clinical verification, such as ST-segment elevation measurements. A second
array is standardised channel-wise for neural processing:

\begin{equation}
  \widetilde{V}_{c}[n]
  =\frac{V_c[n]-\mu_c}{\sigma_c+\varepsilon},
  \label{eq:zscore}
\end{equation}

where $\mu_c$ and $\sigma_c$ are calculated from the valid samples of channel $c$, and
$\varepsilon$ prevents division by zero.

\subsubsection{Operation P4: Zero-Phase Digital Filtering}

Surface ECG recordings may contain respiratory baseline wander, power-line interference,
and high-frequency electromyographic noise. The \texttt{SignalFilter} applies a
three-stage cascade:

\begin{enumerate}[label=\arabic*.,leftmargin=*,align=left]
  \item a third-order Butterworth high-pass filter with a 0.5~Hz cut-off for baseline-wander
  suppression;

  \item second-order infinite-impulse-response notch filters at 50~Hz and 60~Hz with
  quality factor $Q=30$, supporting recordings acquired under either mains-frequency
  environment; and

  \item a fourth-order Butterworth low-pass filter with a 100~Hz cut-off for suppression
  of high-frequency noise.
\end{enumerate}

Every stage uses bidirectional filtering through \texttt{scipy.signal.filtfilt}. The
forward and reverse passes cancel the net phase response, producing zero group delay. This
is necessary to prevent artificial displacement of R peaks, QRS boundaries, ST-segment
onsets, and other diagnostically relevant timings. Filter coefficients and the effective
sampling rate are stored in the record manifest.

\subsubsection{Operation P5: Quality Screening and Integrity Audit}

The \texttt{QualityScreening} component prevents malformed or severely corrupted signals
from entering the foundation model. Four automated gates are applied:

\begin{itemize}
  \item \textbf{Dimensional gate:} verifies the exact shape $12\times5000$;

  \item \textbf{non-finite gate:} rejects any array containing \texttt{NaN}, positive
  infinity, or negative infinity;

  \item \textbf{flatline gate:} flags channels whose voltage standard deviation is below
  $10^{-4}$~mV or which remain effectively constant for more than 1.5 seconds; and

  \item \textbf{saturation gate:} identifies samples at or beyond the $\pm15$~mV
  acquisition limit, indicating amplifier clipping or severe electrode disturbance.
\end{itemize}

The component emits a structured \texttt{QualityAuditResult} containing the overall
pass/fail status, individual gate verdicts, rejection reasons, and per-channel descriptive
statistics. Rejected records are not silently discarded; they remain listed in the manifest
with their exclusion reason.

\subsection{Controlled Lead-Availability Masking}

Signals that pass the integrity audit are supplied to the lead-degradation engine. For each
recording, the engine generates binary masks for the seven availability levels
$\Lambda=(12,10,8,6,4,2,1)$. If $\mathbf{a}_l[c]=0$, channel $c$ is clamped to zero
in the neural-model input:

\begin{equation}
  \widetilde{\mathbf{X}}_l[c,n]
  =\mathbf{a}_l[c]\,\widetilde{\mathbf{X}}[c,n].
  \label{eq:implementation-mask}
\end{equation}

The binary mask is always passed alongside the array, ensuring that zero-valued unavailable
channels are not confused with measured samples whose normalised voltage happens to be
zero. The physical-unit array is masked using the same schedule for independent evidence
verification. Generative inpainting, mathematical lead derivation, and synthetic
reconstruction are excluded. Consequently, all downstream reasoning is performed under
the genuine information scarcity imposed by the selected sensor configuration.

Module 1 stores each output using a compound key containing the patient identifier, record
identifier, preprocessing version, and lead level. The standardised complete recording is
stored once, while degraded variants and masks are generated deterministically from the
manifest to avoid inconsistent copies.

\section{Module 2: Core Reasoning Engine}

Module 2 implements the principal representation-learning and clinical-reasoning
components. It converts the masked multi-lead signal into a dual representation containing
both continuous latent features and discrete, rule-grounded physiological observations.
Candidate diagnoses are then evaluated under the ECG-only control condition and the
knowledge-assisted experimental condition.

The implementation contains three integrated subsystems. These perform foundation-model
adaptation, deterministic extraction of clinical observations, and integration of the
ontological Clinical Knowledge Prior.

\subsection{Foundation-Model Backbone and Residual Bottleneck Adapter}

The software supports two interchangeable waveform encoders. A one-dimensional residual
convolutional network provides a trainable baseline with three stages and channel widths
increasing from 64 to 512. ECG-FM provides the principal pre-trained backbone and uses a
12-layer transformer encoder with 16 attention heads and a feed-forward width of 3072.
Its front-end signal projection is configured with 256 output features, 12 input leads, and
a temporal kernel width of 2.

The backbone remains frozen during the degradation experiment. A compact residual
bottleneck multilayer perceptron maps its activations into the shared representation
$\mathbf{Z}_{\mathrm{ECG}}\in\mathbb{R}^{768}$. The adapter contains:

\begin{enumerate}[label=\arabic*.,leftmargin=*,align=left]
  \item input layer normalisation;
  \item a linear transformation from 512 to 512 features;
  \item a Gaussian Error Linear Unit activation;
  \item dropout with probability $0.20$;
  \item a linear transformation from 512 to 768 features;
  \item a dimension-compatible residual skip path; and
  \item final layer normalisation.
\end{enumerate}

The adapter contains approximately $0.66$ million trainable parameters, depending on
whether the selected backbone requires an explicit skip projection. Keeping this component
small relative to the backbone reduces the risk that it learns a new lead-specific encoder
for every masking level.

\subsection{Adapter Optimisation and Latent Extraction}

The adapter is optimised on PTB-XL Folds~1--8 using multi-label binary cross-entropy
with logits. The frozen backbone ensures that only the adapter and task head receive
gradient updates. Table~\ref{tab:adapter-config} lists the principal optimisation settings.

\begin{table}[H]
  \centering
  \caption{Residual adapter optimisation configuration}
  \label{tab:adapter-config}
  \begin{tabularx}{0.88\textwidth}{@{}lX@{}}
    \toprule
    \textbf{Parameter} & \textbf{Configuration} \\
    \midrule
    Training records & 17,418 recordings from Folds 1--8 \\
    Validation records & 2,183 recordings from Fold 9 \\
    Optimiser & AdamW \\
    Learning rate & $1\times10^{-4}$ \\
    Weight decay & $1\times10^{-2}$ \\
    Maximum epochs & 20 \\
    Dropout probability & 0.20 \\
    Objective & Multi-label binary cross-entropy with logits \\
    Output dimension & 768 \\
    \bottomrule
  \end{tabularx}
\end{table}

Validation loss and macro-level diagnostic performance are monitored after every epoch.
The selected checkpoint is the epoch with the best Fold~9 validation criterion rather than
the final training epoch. The observed training and validation trajectories stabilised
without evidence of continued training improvement accompanied by validation degradation.

The selected adapter generates embeddings for all 2,198 held-out test recordings at every
lead level. For each level $l$, the output matrix has shape

\begin{equation}
  \mathbf{Z}^{(l)}_{\mathrm{test}}\in\mathbb{R}^{2198\times768}.
\end{equation}

Before serialization, each matrix is checked for the expected dimensions, finite values,
and consistent record ordering. The completed extraction produced no non-finite latent
values. Representations are cached so that the control and CKP-assisted reasoning paths
consume exactly the same $\mathbf{Z}_{\mathrm{ECG}}$ for a given record and mask.

\subsection{Deterministic Clinical-Observation Extraction}

The clinical-observation extractor produces discrete physiological biomarkers from the
calibrated, unmasked channels. This explicit representation mirrors the findings used in
clinical ECG interpretation and makes it possible to verify whether a diagnostic claim is
grounded in a visible waveform.

\subsubsection{Cardiac Rhythm and Rate}

R peaks are detected on available rhythm-bearing leads using filtered physical-unit signals.
The median interval between successive valid R peaks is used to estimate heart rate:

\begin{equation}
  \mathrm{HR}
  =\frac{60f_s}{\operatorname{median}(\Delta n_{\mathrm{RR}})}
  \quad\text{beats/min},
  \label{eq:heart-rate}
\end{equation}

where $f_s=500$~Hz and $\Delta n_{\mathrm{RR}}$ is the R--R interval measured in
samples. Rate findings are categorised as bradycardic below 60 beats/min, normal within
the configured clinical interval, or tachycardic above 100 beats/min. Beat-to-beat
dispersion is retained as a rhythm-regularity feature rather than being collapsed into the
rate label.

\subsubsection{Conduction Intervals}

The extractor measures the PR interval from P-wave onset to QRS onset and determines the
QRS duration using onset and offset estimates aggregated across the active leads. All
intervals are stored in milliseconds together with the leads that contributed to the
measurement and a quality flag indicating whether the required landmarks were detected
reliably.

\subsubsection{ST-Segment Morphology}

For each active lead, the ST amplitude is measured at the J point plus 80~ms relative to
the isoelectric TP baseline:

\begin{equation}
  \Delta V_{\mathrm{ST},c}
  =V_c(t_J+80\ \mathrm{ms})-\overline{V}_{\mathrm{TP},c}.
  \label{eq:st-measurement}
\end{equation}

The value is converted to millimetres using the conventional ECG amplitude scale and is
stored separately for elevation and depression. Regional interpretation is performed only
when the required contiguous leads are active.

\subsubsection{T-Wave Polarity}

T-wave polarity is estimated relative to the local isoelectric baseline. Findings are
grouped into anterior, inferior, and lateral anatomical regions. A regional inversion claim
is emitted only when the required lead group is sufficiently observable and the configured
polarity rule is satisfied.

The extractor enforces a strict observability policy across all biomarkers. If a required
lead is absent, the corresponding observation is recorded as \textsc{Unavailable}; it is
never inferred from another channel. Each observation object contains the measured value,
unit, contributing leads, required leads, validity status, and applicable clinical threshold.

\subsection{Ontological Clinical Knowledge Prior}

The CKP represents diagnostic knowledge using a directed graph derived from a curated
cardiac subset of SNOMED CT. Nodes represent diagnoses, morphological findings,
anatomical regions, and lead groups. Edges encode relations such as
\texttt{finding\_site}, \texttt{has\_finding}, \texttt{requires\_lead}, and
\texttt{reciprocal\_to}. Each concept retains its ontology identifier and human-readable
label so that retrieved knowledge remains traceable.

Clinical criteria are converted into deterministic evidence rules. For example, an inferior
myocardial infarction hypothesis requires compatible ST elevation of at least 1.0~mm in
at least two contiguous inferior leads drawn from Leads~II, III, and aVF. The rule is
evaluated only on active leads; if the required regional evidence is not observable, the
finding is returned as unavailable rather than negative.

\subsection{Dual-Path Reasoning Engine}

A shared reasoning transformer evaluates every candidate hypothesis under two conditions:

\begin{itemize}
  \item \textbf{Control path:} receives the masked ECG representation, deterministic
  observations, candidate hypothesis, and associated evidence schema entries; and

  \item \textbf{experimental path:} receives the same information plus retrieved CKP
  relational tokens relevant to the candidate.
\end{itemize}

The hypothesis set, transformer parameters, score head, and verification rules are shared.
The presence of CKP tokens is the only experimental difference. The transformer produces
a posterior score for each diagnosis and a cross-attention matrix linking evidence items
to candidate hypotheses.
The resulting trace records the selected diagnosis, alternative hypotheses, probability
distribution, supporting and opposing claims, attention values, contributing leads, and
uncertainty. This design permits a paired assessment of whether structured knowledge
improves evidence use without changing the input waveform or diagnostic search space.

\section{Module 3: User Interface, Telemetry, and Output Management}

Module 3 implements the researcher-facing interaction, telemetry, reporting, and
visualisation layer. It converts latent representations, posterior scores, cross-attention
matrices, and grounding-verification decisions into human-readable reasoning traces and
machine-readable audit artefacts. The current interface is designed for experimental
inspection and clinical review rather than autonomous diagnostic deployment.

\subsection{Structured Clinical-Reasoning Trace}

For each recording, lead level, and reasoning condition, the system emits a structured
trace with the fields summarised in Table~\ref{tab:trace-fields}.

\begin{table}[H]
  \centering
  \caption{Principal fields in the structured clinical-reasoning trace}
  \label{tab:trace-fields}
  \begin{tabularx}{\textwidth}{@{}l>{\raggedright\arraybackslash}X@{}}
    \toprule
    \textbf{Field} & \textbf{Description} \\
    \midrule
    Record context & Record identifier, patient-safe identifier, split, lead level,
    active leads, model checkpoint, and reasoning condition \\
    Primary hypothesis & Highest-ranked cardiac diagnosis and calibrated posterior
    probability \\
    Alternatives & Remaining candidate diagnoses and their probabilities \\
    Supporting evidence & Structured physiological claims supporting the hypothesis,
    with contributing leads and attention weights \\
    Opposing evidence & Findings that reduce support for the hypothesis \\
    Missing evidence & Required clinical observations that cannot be evaluated from the
    active sensors \\
    Uncertainty & Predictive uncertainty and information-scarcity indicators \\
    Verification & Claim-level \textsc{Supported}, \textsc{Contradicted}, or
    \textsc{Unavailable} verdicts \\
    \bottomrule
  \end{tabularx}
\end{table}

Separating the final diagnosis from the evidence list prevents a correct label from
concealing an unsupported explanation. Attention values are reported as attribution
signals rather than proof of causality and are always interpreted together with the
deterministic verifier.

\subsection{Telemetry and Grounding-Verification Interface}

The telemetry layer passes every claim to the three-state grounding gate. The verifier
checks the required lead set against the active mask and, where sufficient evidence is
observable, evaluates the corresponding physical-unit waveform criterion.

\begin{itemize}
  \item \textsc{Supported} indicates that the relevant evidence is observable and satisfies
  the claim criterion.

  \item \textsc{Contradicted} indicates that the required evidence is observable but does
  not agree with the claim.

  \item \textsc{Unavailable} indicates that the claim cannot be adjudicated because its
  required sensor channels are absent.
\end{itemize}

If $N_{\mathrm{sup}}$, $N_{\mathrm{con}}$, and $N_{\mathrm{unav}}$ are the claim counts
for a trace and
$N_{\mathrm{all}}=N_{\mathrm{sup}}+N_{\mathrm{con}}+N_{\mathrm{unav}}$, the
interface reports

\begin{align}
  \mathrm{Grounded\ Support\ Ratio}
  &=\frac{N_{\mathrm{sup}}}{N_{\mathrm{all}}},
  \label{eq:support-ratio}\\
  \mathrm{Contradicted\ Ratio}
  &=\frac{N_{\mathrm{con}}}{N_{\mathrm{all}}},
  \label{eq:contradicted-ratio}\\
  \mathrm{Unavailable\ Ratio}
  &=\frac{N_{\mathrm{unav}}}{N_{\mathrm{all}}}.
  \label{eq:unavailable-ratio}
\end{align}

The Unavailable Ratio functions as an empirical indicator of how often the reasoning trace
depends on evidence outside the visible sensor space. It is not interpreted alone as a
hallucination rate because a model may explicitly acknowledge unavailable evidence. The
audit log therefore also records whether an unavailable finding was asserted as fact,
appropriately qualified, or used to support the final diagnosis.

\subsection{Automated Reporting and Multi-Cohort Serialization}

Experimental artefacts are serialised in open, inspectable formats. The reporting engine
generates:

\begin{itemize}
  \item Markdown cohort inventories describing dataset sizes, patient counts, split
  assignments, sampling configurations, quality exclusions, and SCP-ECG label mappings;

  \item JSON reasoning traces containing nested hypotheses, evidence claims, lead
  dependencies, attention values, verification verdicts, and uncertainty fields;

  \item CSV summary tables containing macro-AUC, calibration measures, per-lead latent
  statistics, Grounding Quality, Evidence Observability, True Grounding Ratio, AURDC, and
  RFB; and

  \item run manifests recording software configuration, random seed, model checkpoint,
  preprocessing version, knowledge-graph version, and output checksums.
\end{itemize}

Record-level artefacts retain stable identifiers so that a plotted aggregate can be traced
back to the individual waveform, mask, reasoning trace, and verification result that
contributed to it. Personally identifying clinical information is not written to the
experimental output.

\subsection{Graphical Analytics and Visualisation}

The plotting subsystem is implemented with Matplotlib and Seaborn and produces
publication-ready raster or vector figures. The principal visual outputs are:

\begin{enumerate}[label=\arabic*.,leftmargin=*,align=left]
  \item \textbf{Cohort distribution charts}, including grouped and stacked multi-label
  views of the five PTB-XL diagnostic superclasses across the training, validation, and
  held-out test partitions.

  \item \textbf{Signal-quality summaries}, displaying rejection causes, per-channel
  amplitude distributions, and the frequency of flatline or saturation events.

  \item \textbf{Representation diagnostics}, showing latent norms, cosine divergence from
  the complete 12-lead representation, and finite-value audits at each degradation level.

  \item \textbf{Reasoning degradation curves}, plotting Grounding Quality and True
  Grounding Ratio against 12, 10, 8, 6, 4, 2, and 1 available leads for both the control
  and CKP-assisted conditions. Each figure includes dispersion across records, the
  predefined safety threshold at 0.60, and the RFB for both conditions.

  \item \textbf{Claim-state composition plots}, showing how the proportions of supported,
  contradicted, and unavailable claims change as physiological information is removed.
\end{enumerate}

All plots use the same ordered lead axis and condition labels. Aggregate curves are
generated directly from the serialised record-level results rather than from manually
entered values. This ensures that the figures, tables, and quantitative summaries remain
consistent and can be regenerated from a single experiment manifest.
  

\chapter{Results and Discussion}

\section{Experimental Setup and Evaluation Protocol}

The proposed framework was evaluated using the held-out Fold~10 partition of PTB-XL,
comprising 2,198 clinical recordings from 1,880 unique patients
\cite{ptbxl2020}. Patient-stratified isolation was preserved throughout training,
validation, calibration, and testing. Consequently, no recording belonging to a held-out
patient was available during adapter optimisation or hyperparameter selection.

Every test recording was evaluated at seven progressively degraded lead configurations:
12 leads, 10 leads, 8 leads, 6 leads, 4 leads, 2 leads, and 1 lead. The 12-lead
configuration represents the complete clinical acquisition. The 2-lead configuration
approximates emergency or bedside telemetry using Leads~II and V1, while the single-lead
condition represents ambulatory monitoring using Lead~II. The intermediate configurations
provide a controlled transition between these operating settings.

Two matched reasoning paradigms were evaluated:

\begin{enumerate}[label=\arabic*.,leftmargin=*,align=left]
  \item \textbf{ECG-only control.} The control condition operates on the continuous latent
  representation and deterministic observations extracted from the masked ECG. It receives
  no external ontological context.

  \item \textbf{ECG + CKP experimental condition.} The experimental condition receives
  the same waveform, mask, latent representation, observations, hypothesis set, and
  reasoning transformer, but additionally incorporates Clinical Knowledge Prior tokens
  derived from the curated SNOMED CT graph.
\end{enumerate}

This paired design isolates the influence of the CKP. A performance difference cannot be
attributed to a different encoder, test cohort, masking schedule, hypothesis vocabulary, or
verification rule.

\subsection{Evaluation Metrics}

The evaluation jointly measures correctness and observability:

\begin{itemize}
  \item \textbf{Grounding Quality (GQ)} is the proportion of supported claims among all
  observable claims. It measures whether the reasoning agrees with waveform evidence that
  can be checked.

  \item \textbf{Evidence Observability (EO)} is the proportion of clinically required
  evidence that remains physically visible under the active lead subset.

  \item \textbf{True Grounding Ratio (TGR)} combines the two requirements:

  \begin{equation}
    \mathrm{TGR}_l=\mathrm{GQ}_l\times\mathrm{EO}_l.
    \label{eq:tgr-results}
  \end{equation}

  A high GQ with low EO is therefore not interpreted as strong overall grounding.

  \item \textbf{Area Under the Reasoning Degradation Curve (AURDC)} summarises grounding
  resilience over the complete lead-loss trajectory using the normalised numerical
  integration defined in Equation~\ref{eq:aurdc}.

  \item \textbf{Reasoning Failure Boundary (RFB)} is the lowest tested lead count at which
  GQ remains at or above the predefined clinical-safety threshold
  $\tau_{\mathrm{safety}}=0.60$.
\end{itemize}

The results below are cohort-level descriptive measurements. GQ values at individual lead
levels are reported as pooled claim-level ratios, whereas AURDC is calculated from
complete record-level degradation curves and then averaged. The reported AURDC values
therefore should not be reconstructed by numerically integrating the rounded pooled GQ
values shown in the tables.

\section{Quantitative Results and Degradation Analysis}

\subsection{Evidence Observability and Physical Information Loss}

Table~\ref{tab:evidence-retention} summarises the observation-extraction results across the
seven lead configurations. Evidence retention is normalised to the complete 12-lead
condition and is used as the level-wise EO term in the subsequent TGR analysis.

\begin{table}[H]
  \centering
  \small
  \caption{Observable ST evidence under progressive lead degradation}
  \label{tab:evidence-retention}
  \begin{tabularx}{\textwidth}{@{}c>{\raggedright\arraybackslash}Xrrr@{}}
    \toprule
    \textbf{Leads} & \textbf{Active configuration} &
    \textbf{ST elevation} & \textbf{ST depression} &
    \textbf{EO} \\
    \midrule
    12 & Complete standard 12-lead ECG & 53.69\% & 82.30\% & 100.0\% \\
    10 & Leads III and aVF masked & 52.50\% & 82.07\% & 97.8\% \\
    8 & Derived limb leads removed & 37.58\% & 80.85\% & 70.0\% \\
    6 & I, II, V1, V2, V5, and V6 & 36.35\% & 80.35\% & 67.7\% \\
    4 & I, II, V1, and V5 & 25.89\% & 79.39\% & 48.2\% \\
    2 & Telemetry configuration: II and V1 & 22.47\% & 62.60\% & 41.9\% \\
    1 & Ambulatory configuration: II & 0.86\% & 62.06\% & 1.6\% \\
    \bottomrule
  \end{tabularx}
\end{table}

The results demonstrate that lead removal is a loss of anatomical evidence rather than
merely a reduction in sample count. Evidence retention declines gradually from 100\% at
12 leads to 41.9\% at 2 leads, before collapsing to 1.6\% in the single-lead condition.
The small change from 12 to 10 leads indicates that much of the overall evidence remains
observable; however, the removed inferior leads may still be decisive for a specific
regional diagnosis.

ST-elevation evidence is particularly sensitive to the spatial coverage of the acquisition.
Its measured rate falls from 53.69\% in the complete configuration to 22.47\% with
Leads~II and V1 and only 0.86\% with Lead~II alone. Regional myocardial injury is
normally assessed using compatible changes across contiguous anatomical leads. A
single-lead recording cannot provide the spatial agreement required to substantiate an
inferior, anterior, or lateral regional pattern, even if an isolated ST deviation is visible.

ST-depression evidence declines less abruptly, remaining at 62.06\% in the single-lead
condition. This does not mean that the single lead preserves 62.06\% of complete
diagnostic capability. It indicates only that a depression measurement can often be made
on Lead~II; the anatomical distribution, reciprocal relationships, and alternative
explanations remain largely unobservable. The EO column captures this distinction and
prevents an individual measurable feature from being mistaken for complete diagnostic
evidence.

\subsection{Grounding Quality and True Grounding}

Table~\ref{tab:grounding-results} compares the two reasoning conditions. TGR is calculated
from the reported GQ and EO values using Equation~\ref{eq:tgr-results}.

\begin{table}[H]
  \centering
  \small
  \caption{Grounding degradation of the control and CKP-assisted conditions}
  \label{tab:grounding-results}
  \begin{tabular}{@{}crrrrr@{}}
    \toprule
    \textbf{Leads} &
    \textbf{GQ Control} &
    \textbf{GQ CKP} &
    \textbf{Margin} &
    \textbf{TGR Control} &
    \textbf{TGR CKP} \\
    \midrule
    12 & 0.88 & 0.91 & +0.03 & 0.880 & 0.910 \\
    10 & 0.80 & 0.86 & +0.06 & 0.782 & 0.841 \\
    8  & 0.72 & 0.77 & +0.05 & 0.504 & 0.539 \\
    6  & 0.62 & 0.68 & +0.06 & 0.420 & 0.460 \\
    4  & 0.60 & 0.64 & +0.04 & 0.289 & 0.308 \\
    2  & 0.42 & 0.60 & +0.18 & 0.176 & 0.251 \\
    1  & 0.12 & 0.22 & +0.10 & 0.002 & 0.004 \\
    \bottomrule
  \end{tabular}
\end{table}

Both systems degrade as leads are removed, but the CKP-assisted condition achieves a
higher GQ at every operating point. The advantage is modest under complete information
(+0.03 at 12 leads) and becomes largest in the 2-lead configuration (+0.18). This pattern
supports the intended role of the prior: structured knowledge offers limited additional
value when the waveform already provides broad anatomical coverage, but becomes more
useful when the reasoning component must organise a smaller set of observable findings.

The control condition meets the 0.60 threshold at 4 leads and falls below it at 2 leads.
Its RFB is therefore 4 leads. The CKP-assisted condition remains at the threshold with
2 leads and fails at 1 lead, producing an RFB of 2 leads. The knowledge prior consequently
reduces the minimum acceptable lead requirement by two physical leads, corresponding to
one additional tested degradation configuration beyond the control boundary.

TGR provides a more conservative interpretation than GQ alone. At 2 leads, the
CKP-assisted model reaches the GQ threshold of 0.60, but EO is only 0.419, producing a
TGR of 0.251. The reasoning that remains observable is comparatively consistent, yet a
large fraction of the evidence required for complete diagnostic justification is absent.
At one lead, TGR is effectively zero for both systems. Structured knowledge improves the
consistency of evaluable claims but cannot restore the missing anatomical projections.

\subsection{Aggregate Benchmark Measures}

\begin{table}[H]
  \centering
  \caption{Aggregate reasoning-degradation benchmark}
  \label{tab:aggregate-results}
  \begin{tabular}{@{}lrrr@{}}
    \toprule
    \textbf{Metric} & \textbf{ECG-only} & \textbf{ECG + CKP} &
    \textbf{Improvement} \\
    \midrule
    AURDC & 0.70 & 0.81 & +0.11 absolute \\
    Relative AURDC change & --- & --- & +15.7\% \\
    RFB ($\tau=0.60$) & 4 leads & 2 leads & 2 fewer leads required \\
    GQ at 2 leads & 0.42 & 0.60 & +0.18 \\
    GQ at 1 lead & 0.12 & 0.22 & +0.10 \\
    \bottomrule
  \end{tabular}
\end{table}

The CKP-assisted condition achieves an AURDC of 0.81, exceeding the predefined success
criterion of 0.80. This represents an absolute improvement of 0.11 over the control
AURDC of 0.70 and a relative increase of 15.7\%. The RFB result also satisfies the
specified two-lead reduction in the minimum acceptable acquisition. Together, these
measurements show that the CKP improves grounding resilience across the complete
trajectory rather than at only one selected lead level.

The result should nevertheless be interpreted with the TGR findings. An acceptable GQ
does not imply that all evidence required for a definitive diagnosis remains visible.
AURDC and RFB measure the reliability of the reasoning that can be verified, while EO
and TGR reveal the extent of the evidence that has physically survived the degradation.

\section{Comparison with Existing Work}

Established ECG deep-learning studies demonstrate strong diagnostic classification from
complete 12-lead recordings. Ribeiro et al. developed an end-to-end network for six
diagnostic abnormalities \cite{ribeiro2020}, while Strodthoff et al. benchmarked
convolutional and recurrent architectures on PTB-XL using diagnostic and representation
metrics \cite{strodthoff2021}. More recent work such as K-MERL introduces structured
knowledge and a lead-aware encoder for arbitrary-lead representation learning
\cite{kmerl2025}. These studies establish important classification and representation
baselines, but their principal reported outcomes are label- or representation-level
measures rather than deterministic verification of every generated clinical claim against
the physically active leads.

The proposed framework differs in three respects.

\subsection{Evaluation Scope}

Conventional diagnostic metrics reward agreement with the reference label. They do not
determine whether the evidence cited by a reasoning system was observable. The proposed
framework adds a claim-level verification stage that assigns
\textsc{Supported}, \textsc{Contradicted}, or \textsc{Unavailable} to every structured
claim. A correct diagnosis accompanied by evidence from a disconnected lead is therefore
distinguished from a diagnosis supported by the measured waveform.

\subsection{Treatment of Sensor Scarcity}

Reconstruction-based partial-lead approaches attempt to recover missing channels before
classification. Such systems answer a valuable but different question: whether a complete
signal can be estimated from a reduced acquisition. The present framework deliberately
prohibits reconstruction and evaluates the model only on physically observed information.
This exposes the point at which clinical evidence becomes unavailable rather than allowing
synthetic signals to obscure that boundary.

K-MERL is the closest conceptual comparison because it supports arbitrary-lead inputs and
uses structured knowledge. Its evaluation emphasises zero-shot classification and linear
probing performance. The current work instead treats knowledge enhancement as a controlled
reasoning intervention and evaluates whether it improves claim-level grounding, EO,
TGR, AURDC, and RFB.

\subsection{Reduced-Lead Resilience}

At 2 leads, the ECG-only control falls to a GQ of 0.42, while the CKP-assisted model
reaches 0.60. The improvement shows that ontological constraints can preserve consistency
among the claims that remain evaluable under telemetry-level acquisition. However, the
corresponding TGR of 0.251 shows that this condition should not be equated with the
diagnostic completeness of a 12-lead ECG. The contribution is therefore a measurable
extension of grounded reasoning under scarcity, not evidence that reduced-lead monitoring
can replace a complete clinical acquisition.

\section{Research-Based Discussion}

\subsection{The Illusion of Robustness}

A central finding is that label accuracy and reasoning reliability can diverge as sensors are
removed. A model may continue to produce the correct disease label because it has learned
population-level co-occurrences, patient characteristics, or dataset-specific shortcuts.
If its explanation cites a finding from an unavailable lead, the prediction may appear
robust under a conventional classification benchmark even though it is not supported by
the current measurement.

The rapid decline in EO and TGR makes this hidden failure visible. At one lead, the
CKP-assisted model retains a GQ of 0.22, but its TGR is approximately 0.004 because
almost none of the required regional evidence remains observable. The result demonstrates
why GQ, diagnostic accuracy, or confidence should not be interpreted without an explicit
measure of physical observability.

\subsection{Clinical Knowledge as a Reasoning Anchor}

The CKP acts as an inductive constraint within the shared transformer. It links candidate
diagnoses to clinically valid findings, anatomical regions, and lead dependencies. This
reduces the freedom of the reasoning layer to associate a hypothesis with arbitrary latent
features and encourages it to select evidence consistent with the clinical ontology.

The increasing GQ margin at severe degradation supports this interpretation. The largest
gain occurs at 2 leads, where the experimental model improves by 0.18. Nevertheless, the
prior does not generate new sensor evidence. EO is identical for the two conditions, and
both systems fail at the single-lead level. Structured knowledge can organise and constrain
reasoning, but it cannot replace missing physiological measurements.

\subsection{Clinical and Translational Significance}

The 2-lead result is relevant to emergency and bedside telemetry, where Leads~II and V1
are commonly available for rhythm and conduction monitoring. Reaching the GQ threshold
suggests that ontology-guided reasoning may support carefully scoped triage or decision
support when a complete ECG cannot be acquired immediately. The system should expose
its EO, unavailable claims, and uncertainty, and it should request a complete 12-lead
recording whenever the proposed diagnosis depends on regional evidence that the telemetry
configuration cannot observe.

The result does not by itself establish clinical safety. The 2-lead CKP score lies exactly
at the chosen threshold, and the TGR remains low. Prospective validation, clinician review,
external cohorts, and task-specific operating thresholds are required before deployment.

At the single-lead level, both paradigms are below the GQ threshold and have near-zero
TGR. A Lead~II wearable may remain useful for rate and rhythm monitoring, but the
present results do not support using it as the sole basis for complex regional diagnoses
such as localised myocardial infarction. In this condition the appropriate system behaviour
is to abstain, request additional leads, or escalate the case for professional review.

\section{Limitations and Threats to Validity}

Several limitations constrain the interpretation of the current results:

\begin{enumerate}[label=\arabic*.,leftmargin=*,align=left]
  \item The quantitative benchmark is based on one held-out PTB-XL partition. External
  validation on MIMIC-IV-ECG and Chapman-Shaoxing is required to determine whether the
  measured boundaries generalise across institutions, devices, and patient populations.

  \item The 0.60 GQ threshold is a predefined experimental safety criterion rather than a
  universally accepted clinical standard. Threshold sensitivity analysis and clinician-led
  validation are needed for each intended use.

  \item Deterministic observation extraction can itself make measurement errors. The
  verifier reduces model self-consistency bias, but its conclusions remain dependent on
  signal quality, landmark detection, and the completeness of the Evidence Schema.

  \item The reported CKP benefit depends on the coverage and correctness of the curated
  SNOMED CT subgraph. Missing or overly broad relations may reduce the benefit or introduce
  inappropriate constraints.

  \item The supplied aggregate results do not include confidence intervals, repeated-seed
  variation, or statistical significance tests. These are required to establish whether the
  observed margins are stable rather than sample-specific.
\end{enumerate}

These limitations do not invalidate the degradation analysis, but they define the evidence
still required before the measured RFB can be translated into a clinical operating policy.

\section{Chapter Summary}

Progressive lead loss produced a marked decline in both evidence observability and
claim-level grounding. The CKP-assisted model outperformed the ECG-only control at every
lead level, increased AURDC from 0.70 to 0.81, and reduced the RFB from 4 leads to
2 leads. Its largest GQ improvement occurred in the 2-lead telemetry condition. However,
TGR showed that much of the evidence required for complete diagnostic justification was
still unavailable, and both systems failed under the single-lead condition. The principal
conclusion is therefore that structured clinical knowledge can extend the range of
evidence-consistent reasoning, but it cannot substitute for the anatomical information
removed by sensor loss.


%\chapter{Reference}
\begin{thebibliography}{99}
\bibitem{ptbxl2020}
P. Wagner, N. Strodthoff, R.-D. Bousseljot, D. Kreiseler, F. I. Lunze,
W. Samek, and T. Schaeffter, "PTB-XL, a large publicly available
electrocardiography dataset," \textit{Scientific Data}, vol.~7, art.~154, 2020.
DOI: \url{https://doi.org/10.1038/s41597-020-0495-6}.

\bibitem{strodthoff2021}
N. Strodthoff, P. Wagner, T. Schaeffter, and W. Samek,
"Deep learning for ECG analysis: Benchmarks and insights from PTB-XL,"
\textit{IEEE Journal of Biomedical and Health Informatics}, vol.~25, no.~5,
pp.~1519--1528, 2021.
DOI: \url{https://doi.org/10.1109/JBHI.2020.3022989}.

\bibitem{ribeiro2020}
A. H. Ribeiro, M. H. Ribeiro, G. M. M. Paix\~ao, D. M. Oliveira,
P. R. Gomes, J. A. Canazart, M. P. S. Ferreira, C. R. Andersson,
P. W. Macfarlane, W. Meira Jr., T. B. Sch\"on, and A. L. P. Ribeiro,
"Automatic diagnosis of the 12-lead ECG using a deep neural network,"
\textit{Nature Communications}, vol.~11, art.~1760, 2020.
DOI: \url{https://doi.org/10.1038/s41467-020-15432-4}.

\bibitem{kmerl2025}
C. Liu, C. Ouyang, Z. Wan, H. Wang, W. Bai, and R. Arcucci,
"Knowledge-enhanced multimodal ECG representation learning with arbitrary-lead
inputs," in \textit{Findings of the Association for Computational Linguistics:
EMNLP 2025}, pp.~7298--7316, 2025.

\bibitem{samplepaper} A. Author and B. Author, "Title of the research paper", \textit{Journal or Conference Name}, vol.~1, no.~1, pp.~1--10, 2025.
DOI: \url{https://doi.org/10.3844/jcssp.2021.610.623}.
\bibitem{r1} Wu, F., Souza, A., Zhang, T., Fifty, C., Yu, T., and Weinberger, K, "Simplifying graph convolutional networks", \textit{In the International Conference on Machine Learning"}, pages 6861–6871. Available from \url{https://proceedings.mlr.press/v97/wu19e.html}
\bibitem{r2} Zhang, Z. and Fink, O, "Algorithm-informed graph neural networks for leakage detection and localization in water distribution networks", \textit{Reliability Engineering \& System Safety}, page 111494. DOI: \url{https://doi.org/10.1016/j.ress.2025.111494}
\bibitem{r3}, Zhao, F., Liu, J., Zhang, S., Tao, L., Wang, J., Zhang, Z., and Li, N, "A novel method for detecting the leakage duration of water pipe network based on the anchor box mechanism", \textit{In 2025 IEEE 2nd International Conference on Energy and Electrical Engineering (EEE)}, pages 1–5. IEEE, DOI: \url{https://doi.org/10.1109/eee64897.2025.11162744}
\bibitem{samplebook} C. Author, \textit{Title of the Book}, 2nd ed. City, Country: Publisher, 2024.
\bibitem{samplewebsite} Organization Name, ``Title of the webpage,'' 2026. [Online]. Available: \url{https://example.com}. [Accessed: 7-Aug-2026].
\end{thebibliography}

\end{document}